\section{传感器把连续世界变成带时间的数据}

传感器路径不是“读一个寄存器”这么简单。温度、压力、加速度或光强先被敏感元件转换为电压、电流、电容或电阻变化，模拟前端完成偏置、放大与滤波，模数转换器（Analog-to-Digital Converter, ADC）在采样时刻量化输入，数字接口再把样本、状态和时间信息交给控制器。驱动最终看到的整数同时受量程、分辨率、参考电压、校准系数和采样时刻约束。

若一个 12-bit ADC 的参考范围为 0--3.3 V，理想最低有效位约为 $3.3/4096\approx0.806$ mV。读数 2048 只说明输入接近中点；在换算物理量前还要应用零点偏移、增益误差和传感器传递函数。分辨率也不等于精度：噪声、参考源误差、非线性和板级耦合都可能使多个相邻码无法区分真实输入。

采样率必须覆盖信号带宽并留出模拟抗混叠滤波的过渡带。高于 Nyquist 频率的输入会折叠到低频，软件事后无法从样本中判断它原来位于哪个频段。提高数字采样率可以放宽模拟滤波器设计，却会增加总线、FIFO、DMA 和处理器预算；因此采样率是从模拟边界一直影响到内存带宽的系统参数。

\section{I2C、SPI 与采样 FIFO 的边界}

低速配置常用 I2C：控制器按地址选择目标，通过寄存器号读写量程、采样率和中断阈值；目标可用 NACK 表示未接受，也可通过 clock stretching 延长某一位时间。较高速的样本流常用 SPI：片选界定目标和帧，时钟极性/相位决定采样边沿，全双工移位意味着“读寄存器”往往也要同时发送命令或空字节。两者都只保证接口上的 bit/byte 转移，不自动说明样本对应哪个物理时刻。

周期采样设备通常在内部定时器到期时锁存 ADC 结果并写入 FIFO。FIFO 水位达到阈值后产生中断，或由 DMA 批量搬入内存环形缓冲。控制器读取 FIFO 的动作晚于真实采样时刻；若应用需要跨设备对齐，时间戳应尽量靠近采样触发点生成，并说明属于传感器时钟、控制器时钟还是系统时钟。仅用中断到达 CPU 的时间作为样本时间，会把总线等待、中断合并和调度抖动一起计入。

\begin{longtable}{L{0.18\linewidth}L{0.30\linewidth}L{0.32\linewidth}L{0.12\linewidth}}
\toprule
边界 & 保存的状态 & 完成条件 & 常见故障 \\
\midrule
模拟采样 & 量程、参考、滤波与采样沿 & sample/hold 已锁定输入 & 饱和、混叠、参考噪声 \\
ADC 转换 & 通道号、转换序号、校准参数 & 数字码与状态位有效 & 未等待转换完成 \\
串行接口 & 设备地址/片选、寄存器、帧长度 & ACK 或规定数量 bit 已转移 & 模式、端序或帧边界错误 \\
设备 FIFO & 读写指针、水位、溢出标志 & 样本被控制器取走 & 生产速度长期高于消费速度 \\
DMA 环 & 描述符、所有权、时间戳范围 & completion 归还整批缓冲 & CPU 过早复用或 Cache 不可见 \\
应用队列 & 单位、校准版本、序号与时间基 & 消费者确认处理或显式丢弃 & 把旧校准用于新量程 \\
\bottomrule
\end{longtable}

\section{过载、丢样与重新同步}

缓冲区只能吸收短时速率差，不能修复长期过载。采样器以 10 ksample/s 产生 16-bit 三轴数据，未计头部时原始速率为 $10{,}000\times6=60$ kB/s；若控制器最长可能连续 40 ms 得不到总线服务，仅保存这段突发就需要至少 2400 B，再加水位中断延迟和 DMA 描述符切换裕量。队列满时必须定义策略：停止采样、覆盖最旧样本、丢弃最新样本或降低采样率，并通过序号间隙和溢出位让软件知道数据不连续。

错误恢复还要重新建立帧和时间关系。I2C 控制器可能因目标在一字节中途复位而看到总线被拉低；恢复流程可以产生规定数量时钟、发送 STOP 并重新初始化设备。SPI 没有地址和 ACK，片选异常或丢失一个时钟后，主从可能对字边界理解不同，通常要拉高片选结束帧并从已知命令重新开始。设备复位后，旧 FIFO、量程和校准状态都不能沿用；驱动应递增 generation，丢弃无法确定配置归属的残留样本。

从传感器到进程的完整 trace 因此是：物理量在采样沿被锁定，ADC 形成带状态的数字码，FIFO 保存顺序，接口控制器按帧取走，DMA 写入有所有权标记的内存，completion/中断通知驱动，驱动应用单位与校准并发布给进程。任何一层只要丢失序号、时间基或配置代际，数值即使“看起来合理”也可能不再代表正确的物理事件。

\section*{章末练习}

\begin{longtable}{L{0.18\linewidth}L{0.42\linewidth}L{0.30\linewidth}}
\toprule
任务 & 要完成的产物 & 检查要点 \\
\midrule
量化计算 & 计算 12-bit、3.3 V ADC 的理想最低有效位 & 分辨率不等于精度 \\
采样边界 & 说明抗混叠必须在 ADC 之前完成的原因 & 混叠后频率身份不可恢复 \\
时间戳 & 比较采样沿时间和 CPU 中断时间 & 后者包含总线与调度抖动 \\
缓冲预算 & 计算正文三轴样本在 40 ms 停顿中的最低字节数 & 还要为控制延迟留裕量 \\
接口恢复 & 分别给出 I2C 总线卡住和 SPI 帧失步的重新同步入口 & 从已知帧边界重新开始 \\
代际隔离 & 说明传感器复位后为何要丢弃旧 FIFO 样本 & 配置与校准归属可能已改变 \\
\bottomrule
\end{longtable}

\section*{参考解答与要点}

\begin{longtable}{L{0.18\linewidth}L{0.46\linewidth}L{0.26\linewidth}}
\toprule
任务 & 参考解答 & 必须保留的检查点 \\
\midrule
量化计算 & $3.3/4096\approx0.806$ mV/LSB；实际精度还受参考、噪声、偏移与非线性影响。 & 不把码宽直接当准确度 \\
采样边界 & 高频分量采样后会折叠到基带，与真实低频产生相同样本序列；数字处理无法区分二者。 & 模拟滤波位于采样之前 \\
时间戳 & 采样沿代表物理量被锁定的时刻；中断时间还包含转换、FIFO、水位、总线、合并和调度延迟。 & 时间戳要注明时钟域 \\
缓冲预算 & 每样本 6 B，10 ksample/s 在 0.04 s 内产生 $6\times10{,}000\times0.04=2400$ B。 & 这是未计裕量的下限 \\
接口恢复 & I2C 尝试释放时钟并发 STOP 后重配；SPI 拉高片选结束残帧，再从已知命令重新开始。 & 不能从半帧继续解释 \\
代际隔离 & 复位可能恢复默认量程、采样率和校准，残留码无法确定属于哪套配置；generation 变化后应拒绝旧样本。 & 数值合理不代表语义正确 \\
\bottomrule
\end{longtable}
